% ID: SEL-D0038-OB-E01
% ESTADO: cotejada. Segunda pasada visual y borrador compilado revisados.
% Fuente: Selectividad/matematicas_aplicadas/Ua3MATEMÁTICAS APLICADAS examen.pdf
% SHA-256: 4440c15a31a77ee9e052a0668923508e117795b369d569bc02d24b7937366f7a
% Página: 2; opción B; ejercicio original 1.

Se dispone de 160 m de tejido de pana y 240 m de tejido de lana para hacer trajes y abrigos. Se usa 1 m de pana y 2 m de lana para cada traje, y 2 m de pana y 2 m de lana para cada abrigo. Cada traje se vende a 250 \texteuro{} y cada abrigo a 350 \texteuro{}.
\begin{enumerate}[label=\alph*)]
\item \textbf{(2 puntos)} \textquestiondown{}Cuántos trajes y abrigos se deben confeccionar para obtener el máximo beneficio? \textquestiondown{}A cuánto asciende dicho beneficio?
\item \textbf{(0.5 puntos)} \textquestiondown{}Pueden hacerse 60 trajes y 50 abrigos con esas cantidades de tejido? En caso afirmativo, \textquestiondown{}obtendría el máximo beneficio al venderlo todo?
\end{enumerate}
